\section{Methodology, formalization, and numerical evidence}

\emph{The section number \texttt{X} and the equation / theorem / lemma tags
(Lemma X.3.1, Theorem X.4.1, etc.) are placeholders to be assigned
on integration into the full paper.}

This section records the technical and computational content of the
paper: the algebraic identities at the heart of the corrected
\(B_\infty\) framework (§X.3--§X.4), the open challenges that organise
the program forward (§X.4.4, §X.7), the numerical evidence at the
two scales on which our verification operates (§X.5), and the Lean 4
formalisation inventory (§X.6).

The numerical evidence has \textbf{two distinct verified scales}, which
we keep rigorously separate throughout. The \textbf{replication scale}
is \(x = 1.3 \cdot 10^{13}\), at which two independent prime-counting
implementations agree on Koyama's Dominance-of-\(-1\) residue tables
(§X.5.1). The \textbf{analytic-identity scale} is \(K \le 10^{7}\), at
which the corrected \(B_\infty\) identity, the subleading constant
\(C_1\), and the Aoki--Koyama--Mertens drift toward \(e^{-\gamma}\) are
verified across three numerical stacks (§X.5.2--§X.5.4). Numbers
from one scale are not transferred to the other.

\subsection{X.1\quad Notation}

Fix a primitive non-principal Dirichlet character \(\chi\) modulo
\(q \ge 2\), and let \(\rho = \tfrac12 + i\tau\) with \(\tau \neq 0\) be a
simple non-trivial zero of \(L(s,\chi)\). We use:

\begin{itemize}
\item
  \(E_K(\chi,\rho) := \prod_{p \le K} (1 - \chi(p)\,p^{-\rho})^{-1}\)
  (truncated partial Euler product);
\item
  \(c_K(\chi,\rho) := \sum_{n \le K} \mu(n)\,\chi(n)\,n^{-\rho}\)
  (truncated Möbius sum);
\item
  \(D_K(\chi,\rho) := c_K(\chi,\rho)\,E_K(\chi,\rho)\)
  (Dirichlet \(D_K\) statistic);
\item
  \(T_K(\chi,\rho) := \sum_{p \le K}\sum_{k \ge 2} \chi(p)^k / (k\,p^{k\rho})\)
  and its limit \(T_\infty\) with \(B_\infty := \exp(T_\infty)\);
\item
  \(C_1(\chi,\rho) := -L''(\rho,\chi) / (2\,L'(\rho,\chi)^2)\)
  (subleading constant);
\item
  \(\psi\) for the primitive character of conductor \(f \mid q\) inducing
  \(\chi^2\).
\end{itemize}

The branch of \(\log L(2\rho,\psi)\) is fixed by analytic continuation
from \(\mathrm{Re}(s) > 1\), where the absolutely convergent log-Euler
expansion applies, to the boundary line \(\mathrm{Re}(s) = 1\) via
Hadamard--de la Vallée Poussin non-vanishing.

\subsection{X.2\quad Methodology of double verification}

Every numerical claim of §X.5 is computed by \textbf{two independent
implementations in two languages with independent algorithms}.

\begin{itemize}
\item
  \textbf{L1 (primary).} \texttt{mpmath} 1.4 (Python 3.13), 50 decimal places.
  Direct partial-Möbius and partial-Euler evaluation; each zero \(\rho\)
  refined by Muller's method to \(|L(\rho,\chi)| < 10^{-50}\).
\item
  \textbf{L1b (in-language cross-check).} Same library, independent
  algorithm: Hurwitz-zeta expansion
  \(L(s,\chi) = q^{-s}\sum_{a=1}^{q}\chi(a)\,\zeta(s, a/q)\) in place
  of \texttt{mpmath.dirichlet}, central-difference numerical derivatives at
  three step sizes, and an independent linear sieve for \(\mu(n)\).
\item
  \textbf{L2 (cross-language).} PARI/GP 2.17.3 (C), 57 dps default; with
  python-flint 0.8.0 / Arb (FLINT 3.3) at 250 bits as a third stack
  on the worst-case pair.
\end{itemize}

Acceptance gates: agreement to \(\ge 12\) significant digits on \(\rho\)
and to \(\ge 8\) digits on \(L'\), \(L''\), \(C_1\); for residual quantities
where cancellation occurs (the \(B_\infty\) residual), agreement to
\(10^{-12}\) on each of the four component sums separately. Branch
choices and external citation provenance are recorded independently
in L1 and L2.

For the Phase-1 prime-residue replication of §X.5.1, two analogous
stacks \textbf{P1a} and \textbf{P1b} are used: \texttt{primesieve} 12.13 (segmented
wheel sieve) and a hand-rolled plain-C segmented Eratosthenes sieve
with no external dependency.

\subsection{X.3\quad The local Perron double-pole residue}

\textbf{Lemma X.3.1.} \emph{Let \(\chi\) be a primitive non-principal Dirichlet
character and let \(\rho\) be a simple zero of \(L(s,\chi)\). Then for
any \(K > 1\),}
\begin{equation}
\label{eq:res}
\mathop{\mathrm{Res}}_{w = 0}\!\left[\,\frac{K^{w}}{w\,L(w + \rho,\chi)}\,\right]
\;=\;
\frac{\log K}{L'(\rho,\chi)} \;+\; C_1(\chi,\rho).
\end{equation}

\emph{Proof.} Taylor-expand \(L(w + \rho,\chi)\) at \(w = 0\); invert to get
\(1/L(w+\rho,\chi) = (L'(\rho,\chi)\,w)^{-1} - L''(\rho,\chi)/(2\,L'(\rho,\chi)^2) + O(w)\).
Multiply by \(K^w/w = w^{-1} + \log K + O(w)\) and read off the
coefficient of \(w^{-1}\). \(\square\)

Lemma X.3.1 is unconditional given simplicity of \(\rho\). The full
algebraic derivation is in Appendix B §B.2; the identity is also
machine-verified in Lean 4 / Mathlib v4.28.0 (\texttt{LocalPerronResidue.lean},
0 \texttt{sorry}; see §X.6).

\subsection{X.4\quad Identities}

\subsubsection{\texorpdfstring{X.4.1 Corrected \(B_\infty\) identity (unconditional)}{X.4.1 Corrected B\_\textbackslash infty identity (unconditional)}}\label{x.4.1-corrected-b_infty-identity-unconditional}

Our main algebraic result is the following unconditional identity
for the prime-power tail \(T_\infty\).

\textbf{Theorem X.4.1.} \emph{Let \(\chi\) be a primitive non-principal Dirichlet
character of conductor \(q\), let \(\rho\) be a simple zero of
\(L(s,\chi)\) on the critical line, and let \(\psi\) be the primitive
character of conductor \(f \mid q\) inducing \(\chi^2\). Then}
\begin{equation}
\label{eq:Binfty}
T_\infty(\chi,\rho)
\;=\;
\tfrac12 \log L(2\rho,\psi)
\;+\; \mathrm{BPC}_1(\chi,\rho)
\;+\; \mathrm{BPC}_2(\chi,\rho)
\;+\; T_{\ge 3}(\chi,\rho),
\end{equation}
\emph{where}
\[
\mathrm{BPC}_1 = \tfrac12 \sum_{p \mid q,\, p \nmid f}
\log\bigl(1 - \psi(p)\,p^{-2\rho}\bigr),
\]
\[
\mathrm{BPC}_2 = -\tfrac12 \sum_{k \ge 2}\frac1k \sum_p \frac{\chi(p)^{2k}}{p^{2k\rho}},
\qquad
T_{\ge 3} = \sum_{k \ge 3}\frac1k \sum_p \frac{\chi(p)^k}{p^{k\rho}}.
\]
\emph{The four right-hand-side terms are individually finite. The identity
is unconditional given simplicity of \(\rho\).}

The \(k = 1\) prime sum \(\sum_p \chi^2(p) / p^{2\rho}\) is only
conditionally convergent on \(\mathrm{Re}(s) = 1\); its convergence is
supplied by Akatsuka (2013, Lemma 2.1 / eq. (2.5)), which is itself
unconditional (derived from PNT with explicit error term). \(\mathrm{BPC}_2\)
and \(T_{\ge 3}\) are absolutely convergent (minimum exponents
\(\mathrm{Re}(2k\rho) = 2\) and \(\mathrm{Re}(k\rho) = \tfrac32\)). The
full proof is given in Appendix A.

\subsubsection{\texorpdfstring{X.4.2 Subleading constant \(C_1\) and the partial Möbius identity}{X.4.2 Subleading constant C\_1 and the partial Möbius identity}}\label{x.4.2-subleading-constant-c_1-and-the-partial-muxf6bius-identity}

\textbf{Theorem X.4.2.} \emph{Under the hypotheses of Theorem X.4.1, and
assuming every off-target nontrivial zero \(\rho'\) of \(L(s,\chi)\)
with \(|\gamma' - \tau| \le T_K\) is simple (for a zero-avoiding
height \(T_K\) as in the truncated explicit formula of
Inoue 2021 Theorem 1),}
\begin{equation}
\label{eq:cK}
c_K(\chi,\rho) \;=\; \frac{\log K}{L'(\rho,\chi)} \;+\; C_1(\chi,\rho) \;+\; o(1)
\qquad (K \to \infty).
\end{equation}
\emph{The identity is unconditional in \(\rho\) given off-target simplicity.
The rate \(o(1) = O(K^{-1/2+\epsilon})\) is RH-conditional; the
unconditional Soundararajan (2009) bound gives
\(o(1) = O(K^{-1/2}\exp((\log K)^{1/2}(\log\log K)^{14}))\).}

The proof combines Inoue (2021, Theorem 1)'s truncated explicit
formula for \(M^*(K,\chi)\) with Lemma X.3.1 to extract the double-pole
residue at \(\rho\); off-target zero contributions are bounded by
Soundararajan's argument. See Appendix B for the full proof.

\subsubsection{X.4.3\quad The Aoki--Koyama--Mertens constant (Hypothesis AK)}

Aoki--Koyama (2023, \emph{J. Number Theory} \textbf{245}, eq. (1.4), p.~235)
states for a non-principal Dirichlet \(\chi\) that, under DRH,
\[
\lim_{x \to \infty}\Bigl((\log x)^{m}\!\!\prod_{p \le x}(1 - \chi(p)/p^s)^{-1}\Bigr)
\;=\;
\frac{L^{(m)}(s,\chi)}{e^{m\gamma}\,m!} \cdot
\begin{cases}\sqrt 2, & \chi^2 = 1,\ s = \tfrac12,\\ 1, & \text{otherwise,}\end{cases}
\]
with \(m = m_\chi = \mathrm{ord}_{s = 1/2}\,L(s,\chi)\). Specialised to
the regime of this paper (\(m = 1\), \(\rho \ne \tfrac12\), branch
multiplier \(1\)), this reads:
\begin{equation}
\label{eq:AK}
\lim_{K \to \infty} E_K(\chi,\rho)\,\log K \;=\; \frac{L'(\rho,\chi)}{e^{\gamma}}.
\tag*{(AK)}
\end{equation}
This \textbf{corrects} the earlier target \(L'(\rho,\chi) / \zeta(2)\): the
ratio is \(\zeta(2)/e^\gamma \approx 1.0828\). (AK) is DRH-conditional
in the form stated by Aoki--Koyama.

\subsubsection{X.4.4\quad The conditional NDC limit (open)}

Composing (\ref{eq:cK}) with (AK) formally gives
\begin{equation}
\label{eq:NDC}
D_K(\chi,\rho) = c_K(\chi,\rho)\,E_K(\chi,\rho)
\;\longrightarrow\; e^{-\gamma}
\qquad (K \to \infty),
\tag*{(NDC)}
\end{equation}
the corrected Numerical Duality Constant. The composition is
mechanical provided (\ref{eq:cK}) is strengthened to the \textbf{shifted
Perron leading theorem}:
\begin{equation}
\label{eq:Perron-leading}
c_K(\chi,\rho) \;=\; \frac{\log K}{L'(\rho,\chi)} \;+\; o(\log K)
\qquad (K \to \infty).
\tag*{(SP-L)}
\end{equation}
The obstruction is the off-target nontrivial-zero residue aggregate:
even under DRH, an off-target zero \(\lambda\) of multiplicity \(m \ge 2\)
contributes
\[
K^{\lambda - \rho}\,(\log K)^{m-1}\big/\bigl((m-1)!\,(\lambda - \rho)\,a_m(\lambda)\bigr),
\]
and the \((\log K)^{m-1}\) factor is not absorbed by simplicity of
\(\rho\). A literature search (Inoue 2021, Soundararajan 2009,
Ng 2004, Akatsuka 2013) did not surface a theorem that closes
(SP-L); we state (NDC) as \textbf{conditional on (AK) and (SP-L)}
(Question Q:Perron, §X.7).

\subsection{X.5\quad Numerical findings}

\subsubsection{\texorpdfstring{X.5.1 Phase-1 Dominance-of-\(-1\) replication, \(x = 1.3 \cdot 10^{13}\)}{X.5.1 Phase-1 Dominance-of-\/-1 replication, x = 1.3 \textbackslash cdot 10\^{}\{13\}}}\label{x.5.1-phase-1-dominance-of--1-replication-x-1.3-cdot-1013}

We independently replicate the prime-residue counts \(\pi(x; N, a)\)
of Koyama's \emph{nontriv.pdf} for \(N \in \{7, 8, 11, 19, 23\}\) and
\(x \in \{10^{12}, 1.3\cdot 10^{12}, 10^{13}, 1.3\cdot 10^{13}\}\) via
two independent prime-enumeration implementations (\texttt{primesieve} 12.13
and a hand-rolled segmented C sieve). Headline numbers:

\begin{itemize}
\item
  \(\pi(1.3 \cdot 10^{13}) = 445{,}831{,}610{,}611\), cross-checked
  against \texttt{primesieve\ -\/-count} standalone.
\item
  Library-independence: at every one of the four checkpoints, the
  \texttt{primesieve} counts and the hand-rolled C-sieve counts agree on
  every residue class for every \(N\).
\item
  Hardware-independence: a second M1-class machine agrees through
  \(x = 1.3 \cdot 10^{12}\) on every residue class for every \(N\).
\item
  Koyama's identity (3.1), a Dirichlet-orthogonality cross-check on
  the residue-count vector, is verified directly at all \(495\)
  \((N, x, a)\)-cells (worst absolute residual \(1.4 \cdot 10^{-4}\)).
\end{itemize}

Cell-by-cell comparison with Koyama's Tables 3--7 at all four
checkpoints, all moduli:

{\def\LTcaptype{none} % do not increment counter
\begin{longtable}[]{@{}
  >{\raggedright\arraybackslash}p{(\linewidth - 8\tabcolsep) * \real{0.1667}}
  >{\raggedleft\arraybackslash}p{(\linewidth - 8\tabcolsep) * \real{0.2222}}
  >{\raggedleft\arraybackslash}p{(\linewidth - 8\tabcolsep) * \real{0.2222}}
  >{\raggedleft\arraybackslash}p{(\linewidth - 8\tabcolsep) * \real{0.2222}}
  >{\raggedright\arraybackslash}p{(\linewidth - 8\tabcolsep) * \real{0.1667}}@{}}
\toprule\noalign{}
\begin{minipage}[b]{\linewidth}\raggedright
Table
\end{minipage} & \begin{minipage}[b]{\linewidth}\raggedleft
\(N\)
\end{minipage} & \begin{minipage}[b]{\linewidth}\raggedleft
cells
\end{minipage} & \begin{minipage}[b]{\linewidth}\raggedleft
exact
\end{minipage} & \begin{minipage}[b]{\linewidth}\raggedright
substantive disagreement (at \(x = 1.3 \cdot 10^{13}\))
\end{minipage} \\
\midrule\noalign{}
\endhead
\bottomrule\noalign{}
\endlastfoot
3 & 7 & 12 & 11 & 1 cell (\(\Delta = 50\), clean digit-shift profile) \\
4 & 8 & 12 & 1 & 11 (small-\(x\) rows; possible \(x\)-label error in Table 4 draft) \\
5 & 11 & 20 & 19 & 1 cell (\(a = 10\): our \(11{,}503\) vs Koyama \(71{,}711\)) \\
6 & 19 & 18 & 15 & 3 (2 substantive at \(a = 13, 18\); 1 sign flip at small \(x\)) \\
7 & 23 & 30 & 29 & 1 cell (\(\Delta = 100\), clean transposition profile) \\
\textbf{Total} & & \textbf{92} & \textbf{75} & \textbf{17} (74/81 ≈ 91\% excluding the 11 Table-4 small-\(x\) rows) \\
\end{longtable}
}

Comparing to the qualitative dominance-of-\(-1\) statement of Koyama
(\emph{nontriv.pdf}, §3): the signal is \textbf{cleanly reproduced for
\(N = 8\)} (\(-1 = 7\) is the strict largest at \(1.3 \cdot 10^{13}\),
\(\pi(\cdot ; 8, 7) - \pi(\cdot ; 8, 1) = 164{,}958\) vs \(126{,}732\)
and \(102{,}728\)). For \(N = 19\) the ranking of \(-1\) at this checkpoint
is 3rd of 9 non-residues in both our run and Koyama's table ---
consistent with \(-1\) being in the top group but not strictly
dominant. For \(N = 11\) the dominance turns on the disputed cell
\(a = 10\) (Table 5; see below); with Koyama's reported \(71{,}711\),
\(-1\) ranks 2nd of 5 (top group); with our \(11{,}503\), \(-1\) ranks
3rd of 5 (outside top group). Pending reconciliation of that cell.
For \(N = 7\) and \(N = 23\), Koyama himself notes (\emph{nontriv.pdf} p.~19)
that the predicted bias is not yet cleanly observed at this scale
and attributes it to the exceptionally low-lying first zero of the
relevant \(L\)-functions; our data agrees (\(-1\) is 2nd of 3 for
\(N = 7\), mid-rank for \(N = 23\)). Koyama's illustrative estimate
(\emph{nontriv.pdf} p.~19, for \(N = 19\)) places the next sine-wave
peak of the leading complex character at
\(x = e^{33.4} \approx 3.2 \cdot 10^{14}\), indicating that
strict dominance for the harder moduli is expected only at
substantially larger \(x\).

The full replication bundle (source code, build hashes, TSV
outputs, and reproducibility manifest) is deposited as
Supplementary Material S1.

\subsubsection{X.5.2\quad Numerical values of the four Dirichlet pairs}

The four pairs \((\chi, \rho)\) used throughout §X.5.2--§X.5.4 (all
values computed in mpmath at 50 decimal places):

{\def\LTcaptype{none} % do not increment counter
\begin{longtable}[]{@{}
  >{\raggedright\arraybackslash}p{(\linewidth - 8\tabcolsep) * \real{0.2000}}
  >{\raggedright\arraybackslash}p{(\linewidth - 8\tabcolsep) * \real{0.2000}}
  >{\raggedright\arraybackslash}p{(\linewidth - 8\tabcolsep) * \real{0.2000}}
  >{\raggedright\arraybackslash}p{(\linewidth - 8\tabcolsep) * \real{0.2000}}
  >{\raggedright\arraybackslash}p{(\linewidth - 8\tabcolsep) * \real{0.2000}}@{}}
\toprule\noalign{}
\begin{minipage}[b]{\linewidth}\raggedright
Pair
\end{minipage} & \begin{minipage}[b]{\linewidth}\raggedright
\(\chi\) (conductor)
\end{minipage} & \begin{minipage}[b]{\linewidth}\raggedright
\(\rho = \tfrac12 + i\tau\)
\end{minipage} & \begin{minipage}[b]{\linewidth}\raggedright
\(L'(\rho,\chi)\)
\end{minipage} & \begin{minipage}[b]{\linewidth}\raggedright
\(L''(\rho,\chi)\)
\end{minipage} \\
\midrule\noalign{}
\endhead
\bottomrule\noalign{}
\endlastfoot
\(\chi_{-4}/z_1\) & \(\chi_{-4}\) (\(q = 4\)) & \(\tfrac12 + 6.020949 i\) & \(\phantom{-}1.296500 + 0.182765 i\) & \(-1.697050 - 0.554017 i\) \\
\(\chi_{-4}/z_2\) & \(\chi_{-4}\) & \(\tfrac12 + 10.243770 i\) & \(\phantom{-}1.788467 - 0.296776 i\) & \(-3.319767 + 0.755548 i\) \\
\(\chi_5\) & \(\chi_5\) (\(q = 5\)) & \(\tfrac12 + 6.183578 i\) & \(\phantom{-}1.112930 - 0.448830 i\) & \(-1.642973 + 1.035107 i\) \\
\(\chi_{11}\) & \(\chi_{11}\) (\(q = 11\)) & \(\tfrac12 + 3.547041 i\) & \(\phantom{-}1.696582 - 0.250988 i\) & \(-3.121598 + 0.261219 i\) \\
\end{longtable}
}

L1b in-language cross-check (Hurwitz expansion, independent sieve)
agrees with L1 to \(|\Delta L'|, |\Delta L''| \lesssim 6 \cdot 10^{-12}\)
and \(|\Delta C_1| \lesssim 5 \cdot 10^{-13}\). L2 cross-language (PARI/GP
2.17.3) agrees with L1 to \(\ge 11\) decimal digits on every real and
imaginary component. An Arb spot-check at 250 bits on the worst pair
gives interval agreement on \(|L'|\) within \(3 \cdot 10^{-43}\).

\subsubsection{\texorpdfstring{X.5.3 The Aoki--Koyama drift: \(e^{-\gamma}\) vs \(\zeta(2)^{-1}\)}{X.5.3 The Aoki--Koyama drift: e\^{}\{-\textbackslash gamma\} vs \textbackslash zeta(2)\^{}\{-1\}}}\label{x.5.3-the-aokikoyama-drift-e-gamma-vs-zeta2-1}

\emph{(Analytic-identity scale \(K \le 10^{7}\); not transferred from §X.5.1.)}

The modulus \(|D_K|\) statistic at \(K = 2 \cdot 10^{6}\) and \(K = 10^{7}\)
(40 dps, mean over the four pairs):

{\def\LTcaptype{none} % do not increment counter
\begin{longtable}[]{@{}
  >{\raggedright\arraybackslash}p{(\linewidth - 8\tabcolsep) * \real{0.1579}}
  >{\raggedleft\arraybackslash}p{(\linewidth - 8\tabcolsep) * \real{0.2105}}
  >{\raggedleft\arraybackslash}p{(\linewidth - 8\tabcolsep) * \real{0.2105}}
  >{\raggedleft\arraybackslash}p{(\linewidth - 8\tabcolsep) * \real{0.2105}}
  >{\raggedleft\arraybackslash}p{(\linewidth - 8\tabcolsep) * \real{0.2105}}@{}}
\toprule\noalign{}
\begin{minipage}[b]{\linewidth}\raggedright
Quantity
\end{minipage} & \begin{minipage}[b]{\linewidth}\raggedleft
\(K = 2 \cdot 10^{6}\)
\end{minipage} & \begin{minipage}[b]{\linewidth}\raggedleft
\(K = 10^{7}\)
\end{minipage} & \begin{minipage}[b]{\linewidth}\raggedleft
\(\zeta(2)^{-1}\) target
\end{minipage} & \begin{minipage}[b]{\linewidth}\raggedleft
\(e^{-\gamma}\) target
\end{minipage} \\
\midrule\noalign{}
\endhead
\bottomrule\noalign{}
\endlastfoot
Mean \(|D_K| \cdot \zeta(2)\) & \(0.992\) & \(0.974\) & \(1.000\) & \(\zeta(2)\,e^{-\gamma} \approx 0.9237\) \\
Mean \(|E_K \log K|\,e^{\gamma}/|L'|\) & n/a & \(0.942\) & n/a & \(1.000\) \\
\end{longtable}
}

The drift from \(0.992\) to \(0.974\) between \(K = 2 \cdot 10^{6}\) and
\(K = 10^{7}\) is consistent with the AK normalisation \(e^{-\gamma}\)
at the natural \(1/\log K\) finite-size scale and \textbf{incompatible with
the \(\zeta(2)^{-1}\) target} at the same scale. We do not claim
convergence of the complex \(D_K(\chi,\rho)\) from a modulus statistic
alone --- that depends on (SP-L), which is open.

\subsubsection{\texorpdfstring{X.5.4 The \(B_\infty\) identity at the four pairs}{X.5.4 The B\_\textbackslash infty identity at the four pairs}}\label{x.5.4-the-b_infty-identity-at-the-four-pairs}

\emph{(Analytic-identity scale \(K \le 10^{7}\); not transferred from §X.5.1.)}

Identity residual \(|T_K - \mathrm{RHS}|\) for (\ref{eq:Binfty}) at
\(K = 2 \cdot 10^{6}\) (mpmath, 50 dps) and \(K = 10^{7}\) (PARI/GP
2.17.3, closed-form component evaluation):

{\def\LTcaptype{none} % do not increment counter
\begin{longtable}[]{@{}
  >{\raggedright\arraybackslash}p{(\linewidth - 6\tabcolsep) * \real{0.2000}}
  >{\raggedleft\arraybackslash}p{(\linewidth - 6\tabcolsep) * \real{0.2667}}
  >{\raggedleft\arraybackslash}p{(\linewidth - 6\tabcolsep) * \real{0.2667}}
  >{\raggedleft\arraybackslash}p{(\linewidth - 6\tabcolsep) * \real{0.2667}}@{}}
\toprule\noalign{}
\begin{minipage}[b]{\linewidth}\raggedright
Pair
\end{minipage} & \begin{minipage}[b]{\linewidth}\raggedleft
\(K = 2 \cdot 10^{6}\) (L1)
\end{minipage} & \begin{minipage}[b]{\linewidth}\raggedleft
\(K = 2 \cdot 10^{6}\) (L2)
\end{minipage} & \begin{minipage}[b]{\linewidth}\raggedleft
\(K = 10^{7}\) (L2)
\end{minipage} \\
\midrule\noalign{}
\endhead
\bottomrule\noalign{}
\endlastfoot
\(\chi_{-4}/z_1\) & \(2.85 \cdot 10^{-3}\) & \(2.85 \cdot 10^{-3}\) & \(2.58 \cdot 10^{-3}\) \\
\(\chi_{-4}/z_2\) & \(1.66 \cdot 10^{-3}\) & \(1.66 \cdot 10^{-3}\) & \(1.52 \cdot 10^{-3}\) \\
\(\chi_5\) & \(4.24 \cdot 10^{-5}\) & \(4.24 \cdot 10^{-5}\) & \(1.22 \cdot 10^{-5}\) \\
\(\chi_{11}\) & \(3.34 \cdot 10^{-5}\) & \(3.34 \cdot 10^{-5}\) & \(1.75 \cdot 10^{-5}\) \\
\end{longtable}
}

L1 and L2 agree to all displayed digits at \(K = 2 \cdot 10^{6}\)
(stack difference \(\le 10^{-8}\)). For the clean-character pairs
\(\chi_5\), \(\chi_{11}\) (where \(\chi(2) \ne 0\) and there is no bad-prime
contribution to \(\mathrm{BPC}_1\)), the residual decays by a factor
\(1.9\)--\(3.5\) from \(K = 2 \cdot 10^{6}\) to \(K = 10^{7}\), bracketing
the \(\sqrt{5} \approx 2.24\) predicted by the \(K^{-1/2} / \log K\) rate
of the boundary-line conditional tail (Akatsuka 2013 eq. (2.5)). The
\(\chi_{-4}\) pairs show systematically larger residuals consistent
with the bad-prime \(p = 2\) contribution to \(\mathrm{BPC}_1\).

\subsubsection{X.5.5\quad Two negative elliptic-curve findings}

We record two negative findings on the elliptic-curve side, both
deferred to the supplementary record for the design data, source,
and resampling diagnostics:

\begin{itemize}
\item
  \emph{Conductor-confounded rank trend.} A multivariate OLS regression
  of \(E[C_1^2]\) on \((\mathrm{rank}, \log N)\) across 19 weight-\(2\)
  elliptic curves shows a stable \(\log N\) coefficient and an
  \emph{unstable} rank coefficient (bootstrap \(95\%\) CI for the rank
  coefficient includes zero; one rank-\(3\) anchor swings the
  coefficient by \(63\%\)). We describe this as a
  conductor-confounded trend, not a rank law (Question Q:conductor,
  §X.7).
\item
  \emph{Raw \(\mathrm{Sym}^2 / \langle f, f \rangle\) proportionality.}
  The raw ratio ranges over seven orders of magnitude across
  \(\{37a_1, 389a_1, \Delta\}\) and is empirically falsified in its
  raw form (Question Q:Sym2, §X.7).
\end{itemize}

\subsection{X.6\quad Lean 4 / Mathlib4 formalisation}

We accompany the paper with a Lean 4 / Mathlib4 lake project
(\texttt{formal-conjectures/} directory of the supplementary archive).
The toolchain is \texttt{leanprover/lean4:v4.28.0}; Mathlib is pinned at
commit \texttt{8f9d9cff6bd728b17a24e163c9402775d9e6a365}. The Lean
inventory fixes
the \textbf{statements} of every identity of §X.4, ensures normalisations
and branch conventions are syntactically explicit, and records each
statement's proof status against a public audit trail.

\textbf{Build status.} \texttt{lake\ build\ FormalConjectures} succeeds on all
\textbf{9 files} in \texttt{formal-conjectures/} with \textbf{5 \texttt{sorry} warnings},
each annotated in-source as \texttt{MATHLIB-PREREQ:} or \texttt{RESEARCH-OPEN:}.
\textbf{Six files are fully proved (0 \texttt{sorry})}: all
the algebraic content of §X.3, §X.4, the smoothed
explicit-formula chain, the Mertens spectroscope universality
statement, the Farey bridge identity, and DPAC for \(K \le 4\).
No \texttt{axiom} is introduced anywhere in the project. The full
per-\texttt{sorry} inventory is in \texttt{LEAN\_SORRY\_STATUS.md} of the
reproducibility bundle.

\textbf{Status of headline theorems} in §X.3--§X.4:

{\def\LTcaptype{none} % do not increment counter
\begin{longtable}[]{@{}
  >{\raggedright\arraybackslash}p{(\linewidth - 2\tabcolsep) * \real{0.5000}}
  >{\raggedright\arraybackslash}p{(\linewidth - 2\tabcolsep) * \real{0.5000}}@{}}
\toprule\noalign{}
\begin{minipage}[b]{\linewidth}\raggedright
Paper object
\end{minipage} & \begin{minipage}[b]{\linewidth}\raggedright
Lean status
\end{minipage} \\
\midrule\noalign{}
\endhead
\bottomrule\noalign{}
\endlastfoot
Lemma X.3.1 (local Perron residue) & \textbf{Lean-verified unconditional} (\texttt{LocalPerronResidue.lean}, 0 \texttt{sorry}) \\
Theorem X.4.1 (corrected \(B_\infty\) identity) & \textbf{Lean-verified conditional} on the convergence-of-partial-sum hypothesis derived in Appendix A (\texttt{CorrectedBInfty.lean}, 0 \texttt{sorry}) \\
Theorem X.4.2 (\(c_K\) leading + subleading) & Pen-and-paper proof in Appendix B; off-target-zero-simplicity hypothesis stated explicitly \\
(AK), (SP-L), (NDC) & (AK) cited (Aoki--Koyama 2023); (SP-L) and (NDC) open challenges (Q:Perron, §X.7) \\
DPAC & \textbf{Lean-verified for \(K \le 4\) unconditional} (\texttt{DPAC\_closure\_attempt.lean}, 0 \texttt{sorry}); general \(K\) open (LI-class), with obstruction certificate via Pólya 1913 \\
\end{longtable}
}

{\def\LTcaptype{none} % do not increment counter
\begin{longtable}[]{@{}
  >{\raggedright\arraybackslash}p{(\linewidth - 4\tabcolsep) * \real{0.3333}}
  >{\raggedright\arraybackslash}p{(\linewidth - 4\tabcolsep) * \real{0.3333}}
  >{\raggedright\arraybackslash}p{(\linewidth - 4\tabcolsep) * \real{0.3333}}@{}}
\toprule\noalign{}
\begin{minipage}[b]{\linewidth}\raggedright
Paper object
\end{minipage} & \begin{minipage}[b]{\linewidth}\raggedright
Lean file
\end{minipage} & \begin{minipage}[b]{\linewidth}\raggedright
Status
\end{minipage} \\
\midrule\noalign{}
\endhead
\bottomrule\noalign{}
\endlastfoot
Boundary residue \(R_0 = -2\) for a Gaussian-cutoff Mellin-shift explicit formula (companion strand) and its algebraic-glue chain & \texttt{SmoothedDwfFormula\_full.lean} & \textbf{THEOREM (chain), 0 \texttt{sorry}.} All 17 algebraic-glue lemmas closed unconditionally; the two analytic prerequisites \texttt{mellin\_decay} (Stirling on \(\Gamma\) vertical strips) and \texttt{inv\_zeta\_polynomial\_growth} (Titchmarsh §3.11) are now stated as explicit hypotheses on the theorems that consume them, both Mathlib v4.28.0 gaps. \\
Lemma X.3.1 (local Perron residue) & \texttt{LocalPerronResidue.lean} & \textbf{THEOREM (0 \texttt{sorry}).} The residue identity is stated as a \texttt{Tendsto} limit at \(L\) analytic with simple zero at \(0\) (the general-\(\rho\) case reduces by \(L \mapsto L(\cdot + \rho)\)). \\
Theorem X.4.1 (\(B_\infty\) identity) & \texttt{CorrectedBInfty.lean} & \textbf{THEOREM (0 \texttt{sorry}), conditional on \texttt{h\_convergence}.} The four-component identity is proved against \texttt{noncomputable\ def}s of \(T_\infty\), \(T_{\ge 3}\), \(\mathrm{BPC}_1\), \(\mathrm{BPC}_2\), \(L\) given an added hypothesis \texttt{h\_convergence\ :\ Tendsto\ T\_K\ atTop\ (nhds\ (RHS))}. This hypothesis packages exactly the four analytic inputs of the pen-and-paper proof in Appendix A (Akatsuka 2013, log-Euler-product, imprimitive induction, geometric tails); the Lean proof uses \texttt{Classical.epsilon\_spec} + \texttt{tendsto\_nhds\_unique} and is three lines. \\
Farey bridge identity & \texttt{FareyBridgeIdentity.lean} & \textbf{THEOREM (0 \texttt{sorry}), conditional on a Ramanujan-sum decomposition hypothesis} packaging the Hardy--Wright Theorem 304 input (\texttt{c\_q(p)\ =\ \textbackslash{}mu(q)} for \(\gcd(p,q) = 1\)). Two helper lemmas (\texttt{mertens\_eq\_pred\_add\_moebius}, \texttt{sum\_moebius\_Icc\_eq\_mertens}) closed without \texttt{sorry}. \texttt{MATHLIB-PREREQ:\ Mathlib.NumberTheory.RamanujanSum} for the unconditional version. \\
Mertens spectroscope universality & \texttt{MertensSpectroscopeUniversality.lean} & \textbf{THEOREM (0 \texttt{sorry}), conditional on an explicit-formula asymptotic hypothesis} (Soundararajan 2009 Theorem 1 input). Proof: \texttt{Filter.tendsto\_atTop.mpr\ h\_explicit\_formula}. \\
Farey sign pattern & \texttt{FareySignPattern.lean} & \textbf{NEGATIVE.} The pointwise version is falsified at \(p = 237{,}733\) and \(p = 243{,}799\) (recorded as theorems, not axioms --- the project's ``no \texttt{axiom}'' convention is preserved). The density-one surviving version is stated as a \texttt{Tendsto}. 3 \texttt{sorry}s. \\
Dirichlet Polynomial Avoidance (DPAC) --- partial closure + bridges & \texttt{DPAC\_closure\_attempt.lean}, \texttt{DirichletPolynomialAvoidance.lean}, \texttt{DPAC\_full.lean} & \textbf{PARTIAL.} \texttt{DPAC\_closure\_attempt.lean} (0 \texttt{sorry}) proves DPAC unconditionally for \(K \in \{2, 3, 4\}\) using only \(0 < \mathrm{Re}(\rho) < 1\) (\texttt{dpac\_K\_eq\_2}, \texttt{dpac\_K\_eq\_3}, \texttt{dpac\_K\_eq\_4}, \texttt{dpac\_le\_4}). It also reformulates the open case as \texttt{FiniteLogRatioLI} and records the obstruction certificate (Pólya 1913 discreteness of the exponential-polynomial zero set + a single open avoidance statement). The headline conjecture for general \(K\) remains \texttt{sorry} in \texttt{DPAC\_full.lean:297} and \texttt{DirichletPolynomialAvoidance.lean:48}, diagnostically comparable to the Linear Independence Hypothesis for \(\zeta\)-zero ordinates; the four explicit phase-avoidance bridges (\texttt{dpac\_of\_logPrimePhaseAvoidance}, \ldots, \texttt{dpac\_of\_certifiedZetaZeroSample}) are closed without \texttt{sorry}. \\
\end{longtable}
}

The role of the Lean artifact is to fix the statements and provide
a publicly inspectable audit trail of the proof obligations
remaining.

\subsection{X.7\quad Open challenges}

The following structure the next phase of the program.

\begin{quote}
\textbf{Q:Perron (Shifted Perron leading theorem).} Prove (SP-L)
(\ref{eq:Perron-leading}) for primitive non-principal \(\chi\) and
simple non-central \(\rho\).
\end{quote}

Sufficient packages: (a) all off-target zeros simple, and
\(Z_\mathrm{simple}(K, T_K) := \sum_{\rho' \ne \rho,\,|\gamma'| \le T_K} K^{\rho'-\rho}/[(\rho'-\rho)\,L'(\rho',\chi)] = o(\log K)\)
at a zero-avoiding height \(T_K \asymp K(\log K)^{-B}\); (b) a
Dirichlet shifted second moment
\(\sum_{\rho'}^\mathrm{mult} |L(\rho' + \alpha,\chi)|^{-2} \ll_\chi (\log K)^{O(1)}\).
The total-Möbius bounds of Soundararajan type are too coarse to
isolate the pointwise cancellation at the \(\log K\) scale.

A naïve GL(1) transfer of the GL(2) halo bound (cluster-summed
contour residues over a halo region, giving \(\ll T^{7/4+\varepsilon}\)
at \(T\)-scale) yields only \(K^{1/2+\varepsilon}\) for the off-target
residue aggregate at \(K\)-scale, far above the \(o(\log K)\) target ---
the halo route in its present form does not close (SP-L). The
cluster-summed-residue pivot does, however, replace the rooted
Palm wall that obstructs the termwise estimate. Details in the
supplementary record.

\begin{quote}
\textbf{Q:EC-recip (GL(2) reciprocal-derivative control).} Prove a
fixed-curve theorem for
\(\sum_\gamma \widehat W(i\gamma)\,e^{i\gamma u} / L'(E, 1 + i\gamma)\)
giving cancellation \(o(u^r)\), or a minimum-modulus estimate
on a vertical line with explicit exponent \(< 2\).
\end{quote}

Without a GL(2) analogue of Aoki--Koyama, the EC side remains at the
level of quantitative ensemble evidence.

\begin{quote}
\textbf{Q:DPAC (Dirichlet Polynomial Avoidance).} Prove DPAC for
general \(K\). We give an unconditional Lean proof for
\(K \in \{2, 3, 4\}\); the general case reduces, via Pólya 1913
discreteness of the finite-exponential-polynomial zero set, to
a single open avoidance statement at \(\zeta\)-zero ordinates
diagnostically comparable to the Linear Independence Hypothesis.
\end{quote}

\textbf{Further questions} (from the EC and ensemble negatives of §X.5.5;
deferred to the supplementary record):

\begin{itemize}
\item
  \emph{Q:conductor} --- Replicate (\ref{eq:W2}) on a curve set in which
  rank and \(\log N\) are not collinear, to separate the conductor
  contribution from any genuine rank dependence.
\item
  \emph{Q:Sym2} --- Identify a completed / archimedean-corrected
  \(\mathrm{Sym}^2\) normalisation replacing the empirically falsified
  raw \(\mathrm{Sym}^2 / \langle f, f \rangle\) proportionality.
\item
  \emph{Q:EC-NDC} --- Find a normalisation of \(D_K^E\) for which the
  universal limit exists and survives a null-control gate. The
  sharp-cutoff form \(D_K^E \cdot \zeta(2) \to 1\) is falsified
  through \(K = 10^{6}\); smoothed variants pass empirically but
  also pass a null-control gate against predeclared null
  transformations.
\end{itemize}

\subsection{X.8\quad Code, data, and certificate availability}

All scripts, refined zero data, numerical-table CSVs, convergence
logs, the Lean 4 lake project, and the reproducibility manifest will
be deposited as a single self-contained reproducibility bundle
(Supplementary material S1), mirrored at a Zenodo DOI at acceptance.
The bundle pins all software versions: Lean toolchain
\texttt{leanprover/lean4:v4.28.0}, Mathlib commit
\texttt{8f9d9cff6bd728b17a24e163c9402775d9e6a365}, \texttt{mpmath} 1.4,
PARI/GP 2.17.3, FLINT 3.3 / python-flint 0.8.0. The Phase-1
Dominance-of-\(-1\) replication bundle is included as the
Supplementary Material S1 archive. Each numerical
table in §X.5 cites the L1 script and L2 reproducer; each external
theorem cited in §X.4 has its PDF retrieval recipe, page/equation,
and verbatim quote recorded in the citation audit
(Supplementary S2, \emph{Citation audit}).

\subsection{References (section)}

External references cited in §X.3--§X.7. Full page-and-equation
provenance for each is in Supplementary S2 (citation audit).

\begin{itemize}
\item
  \textbf{Akatsuka, H.} (2013). \emph{The Euler product for the Riemann zeta
  function in the critical strip}. The boundary-line Mertens-type
  partial-summation estimate
  \(\sum_{p \le X} \chi(p) / p^{1 + 2i\tau} = c(\chi, \tau) + O(1 / \log X)\)
  used in §X.4.1, Appendix A §A.2.3; Lemma 2.1 / eq. (2.5).
  Unconditional (derived from PNT with explicit error term).
\item
  \textbf{Aoki, T. and Koyama, S.} (2023). \emph{Generalised Mertens constant
  for Dirichlet \(L\)-functions.} J. Number Theory \textbf{245}, 233--262;
  eq. (1.4), p.~235. The \(e^{-m\gamma}\)-normalised partial-Euler-product
  asymptotic at noncentral zeros under DRH. Cited as Hypothesis (AK)
  in §X.4.3.
\item
  \textbf{Davenport, H.} \emph{Multiplicative Number Theory}, 3rd ed.,
  GTM \textbf{74}, Springer 2000. Hadamard--de la Vallée Poussin
  non-vanishing on \(\mathrm{Re}(s) = 1\).
\item
  \textbf{Hardy, G.H. and Wright, E.M.} \emph{An Introduction to the Theory of
  Numbers}, 6th ed., Oxford 2008. Theorem 304 (Ramanujan-sum identity
  at primes).
\item
  \textbf{Inoue, S.} (2021). \emph{Truncated explicit formula for \(M(x, \chi)\).}
  Journal de Théorie des Nombres de Bordeaux \textbf{33}, Theorem 1. Used
  in Appendix B §B.1, §B.3 to set up the contour integration for
  Theorem X.4.2.
\item
  \textbf{Montgomery, H.L. and Vaughan, R.C.} \emph{Multiplicative Number
  Theory I. Classical Theory}, Cambridge 2007. Theorem 9.4 (textbook
  bad-prime correction structure).
\item
  \textbf{Ng, N.} (2004). \emph{The distribution of the summatory function of
  the Möbius function.} Proc. London Math. Soc. \textbf{89}(3), 361--389.
  Discussed alongside Soundararajan 2009 for the off-target zero
  aggregate (Question Q:Perron, §X.7).
\item
  \textbf{Pólya, G.} (1913). \emph{Über die Nullstellen gewisser ganzer
  Funktionen.} Classical discreteness of the zero set of a finite
  exponential polynomial; used in \texttt{DPAC\_closure\_attempt.lean} for
  the obstruction certificate.
\item
  \textbf{Soundararajan, K.} (2009). \emph{Partial sums of the Möbius
  function.} Ann. of Math. \textbf{170}(2), 1409--1422; Theorem 1. The
  RH-conditional and the unconditional rate bounds for the Möbius
  partial sum used in Theorem X.4.2 and Appendix B §B.4.
\item
  \textbf{Tenenbaum, G.} \emph{Introduction to Analytic and Probabilistic
  Number Theory}, 3rd ed., GSM \textbf{163}, AMS 2015. Chapter II.5
  (Hadamard--de la Vallée Poussin treatment).
\item
  \textbf{Titchmarsh, E.C.} \emph{The Theory of the Riemann Zeta-Function},
  2nd ed., Heath-Brown rev., OUP 1986. §3.11 (the polynomial-growth
  bound on \(1/\zeta(s)\) used in \texttt{SmoothedDwfFormula\_full.lean} and
  cited as a Mathlib prerequisite in §X.6).
\end{itemize}
